% ======================================================================
% Paper Section Template
% Use this structure when writing a paper/research note section.
% ======================================================================

% ---------- Introduction ----------
\section{Introduction}
\label{sec:intro}

\paragraph{Motivation.}
% Why does this problem matter? Real-world impact, scientific importance.

\paragraph{Problem.}
% Formal statement: what is the task, what are the inputs/outputs.

\paragraph{Limitations of existing work.}
% What's wrong with current methods? Be specific and cite.

\paragraph{Our contribution.}
% What do we propose? Bullet list:
% \begin{itemize}
%   \item Novel method X that addresses limitation Y
%   \item Theoretical analysis showing Z
%   \item Empirical results demonstrating W improvement
% \end{itemize}

% ---------- Related Work ----------
\section{Related Work}
\label{sec:related}

% Position your work relative to:
% - Classical / traditional methods
% - Learning-based methods
% - Hybrid approaches
% - Specific sub-areas

% Use a table for comparison if helpful:
% \begin{table}[t]
%   \centering
%   \caption{Comparison with prior work.}
%   \label{tab:related}
%   \begin{tabular}{lccc}
%     \toprule
%     Method & Handles X & Real-time & Trainable \\
%     \midrule
%     Method A & Yes & No & No \\
%     Method B & No & Yes & Yes \\
%     \textbf{Ours} & \textbf{Yes} & \textbf{Yes} & \textbf{Yes} \\
%     \bottomrule
%   \end{tabular}
% \end{table}

% ---------- Method ----------
\section{Proposed Method}
\label{sec:method}

\subsection{Problem Formulation}
% Math notation. Define all variables.
% \begin{equation}
%   \min_{\theta} \mathbb{E}_{(\mathbf{x}, \mathbf{y}) \sim \mathcal{D}} [\ell(f_\theta(\mathbf{x}), \mathbf{y})]
% \end{equation}

\subsection{Proposed Approach}
% Detailed method description. Algorithm pseudocode if applicable.

% \begin{algorithm}[t]
%   \caption{Our proposed algorithm}
%   \label{alg:ours}
%   \begin{algorithmic}
%     \State \textbf{Input:} ...
%     \State \textbf{Output:} ...
%     \For{each iteration}
%       \State ...
%     \EndFor
%   \end{algorithmic}
% \end{algorithm}

\subsection{Theoretical Analysis}
% Convergence, complexity, invariance, etc.
% \begin{theorem}
%   Under assumptions A1-A3, Algorithm 1 converges at rate O(1/T).
% \end{theorem}
% \begin{proof}
%   ...
% \end{proof}

% ---------- Experiments ----------
\section{Experiments}
\label{sec:experiments}

\subsection{Experimental Setup}

\paragraph{Datasets.}
% List datasets, sizes, splits, preprocessing.

\paragraph{Baselines.}
% Methods compared, with brief description and citations.

\paragraph{Implementation Details.}
% Architecture, hyperparameters, training, hardware.
% Include a table:
% \begin{table}[t]
%   \centering
%   \caption{Hyperparameters.}
%   \label{tab:hyperparams}
%   \begin{tabular}{ll}
%     \toprule
%     Parameter & Value \\
%     \midrule
%     Learning rate & $10^{-4}$ \\
%     Batch size & 32 \\
%     Optimizer & Adam \\
%     Epochs & 100 \\
%     \bottomrule
%   \end{tabular}
% \end{table}

\subsection{Main Results}
% Table with key comparison:
% \begin{table}[t]
%   \centering
%   \caption{Quantitative results. Best in bold, second underlined.}
%   \label{tab:main-results}
%   \begin{tabular}{lcc}
%     \toprule
%     Method & PSNR $\uparrow$ & SSIM $\uparrow$ \\
%     \midrule
%     Baseline A & 28.3 $\pm$ 0.2 & 0.892 $\pm$ 0.003 \\
%     Baseline B & 29.1 $\pm$ 0.3 & 0.901 $\pm$ 0.004 \\
%     \textbf{Ours} & \textbf{30.7 $\pm$ 0.2} & \textbf{0.934 $\pm$ 0.002} \\
%     \bottomrule
%   \end{tabular}
% \end{table}

\subsection{Ablation Studies}
% What matters and why? Remove each component.
% \begin{figure}[t]
%   \centering
%   \includegraphics[width=\columnwidth]{results/figures/ablation.pdf}
%   \caption{Ablation study: impact of each component.}
%   \label{fig:ablation}
% \end{figure}

% ---------- Discussion ----------
\section{Discussion}
\label{sec:discussion}

\paragraph{Limitations.}
% Be honest. What doesn't your method handle?

\paragraph{Future Work.}
% What would you do next with unlimited time?

% ---------- Conclusion ----------
\section{Conclusion}
\label{sec:conclusion}

% Summary in 3-5 sentences: problem, method, key result, impact.
